% =====================================================================
% Section II: Related Work
% =====================================================================
\section{Related Work}
\label{sec:related}

\subsection{Model Based Reconstruction and Hand Crafted Priors}

Statistical reconstruction of photon limited data begins with the
multiplicative EM iteration of Shepp and Vardi \cite{shepp1982}, whose
likelihood ascent is unregularised and therefore recovers noise as it
converges. The standard remedies add a prior: total variation and its
constrained formulations \cite{sidky2008}, anisotropic diffusion
\cite{perona1990} and its fuzzy variants, and plug and play splittings that
insert a generic denoiser into an ADMM iteration
\cite{venkatakrishnan2013}. Fractional order penalties were introduced to
soften the staircase artefacts of first order total variation and to give a
continuous control over the smoothness class; in CT they have been used both
as a fractional total variation and as an adaptive fractional order penalty
whose exponent is tuned per image \cite{zhang2016fractional}. All of these
methods share a free strength parameter, and the present work shows that its
optimum moves by an order of magnitude between operating points, so a value
carried over from one setting silently degrades another.

Our own starting point \cite{singh2026spl} differs from the above in where
the prior acts: instead of penalising spatial derivatives, it applies a
zero-phase spectral operator that combines a fractional-order envelope with
narrowband stop bands, which makes it possible to attenuate structured
interference without smoothing edges. That letter, however, required both the
order and the notch positions as inputs, which is the limitation removed
here.

\subsection{Learned Reconstruction}

Learned methods for low-dose CT fall into three families. Post-processing
networks map a noisy reconstruction to a clean one and are represented by the
residual encoder decoder of Chen \emph{et al.} \cite{chen2017redcnn} and by
adversarially trained variants \cite{yang2018wgan}; architecture search has
been used to reduce their cost \cite{lu2023m3nas}. Unrolled methods embed the
forward operator in the network, which improves data consistency and reduces
the amount of training data required; learned primal-dual \cite{adler2018lpd}
is the canonical example, with dual domain and optimisation inspired variants
following \cite{wu2021drone,xiang2021fista}. Generative methods learn a prior
over images and use it during inference: diffusion posterior sampling
\cite{chung2023dps} conditions an unconditional score model on the
measurements, and recent CT specific work has addressed the cost and the
generalisation of such priors \cite{gao2024corediff} and their use for
artefact reduction in dual domains \cite{liu2024mar}. Between the trained and
the training free extremes sits the deep image prior
\cite{ulyanov2018dip,baguer2020dip}, which fits an untrained network to a
single measurement and is therefore the natural neural comparator for a
method that uses no training data at all.

\subsection{Robustness, Generalisation and Benchmarking}

The reliability of learned reconstruction has been studied from two
directions. Antun \emph{et al.} \cite{antun2020} showed that small
perturbations can produce severe artefacts and that clinically relevant
structures can be missed, and later analyses examined how far such
instabilities extend in practice \cite{genzel2023}. Independently, controlled
benchmarks have questioned how much of the reported progress survives a
common training and evaluation protocol: Eulig \emph{et al.} \cite{eulig2024}
retrained a range of low-dose denoisers under identical conditions and found
that newer architectures do not consistently beat the older residual encoder
decoder. Domain shift is the practical face of the same issue, and it
motivates the generalisation oriented design of recent diffusion models for
CT \cite{gao2024corediff}. The present work contributes a controlled
measurement of that effect across a full method jury: the same trained
weights are evaluated on a second public dataset without adaptation, and the
loss is reported per family.

\subsection{Parameter Selection}

Choosing the regularisation strength from the data is classical in inverse
problems, through the discrepancy principle, the L curve or unbiased risk
estimation, but it is comparatively rare in CT practice, where the strength
is usually fixed by protocol or by a sweep against a reference. For Poisson
data the natural criterion is the deviance, which has a known expectation
under the model and therefore provides a calibrated stopping rule. We use it
to select the fractional order, and we show that a bounded candidate set,
which is the usual implementation choice, turns the criterion into a fixed
choice whenever the data demand stronger regularisation than the largest
candidate.

\subsection{Structured Artefacts}

Ring and stripe artefacts arise from detector gain drift, defective channels
and scintillator inhomogeneity. Sinogram domain corrections estimate a per
channel gain, while image domain methods work in polar coordinates or in a
transform domain in which rings become oriented structures. Our notch rule
belongs to the image domain family and is explicitly limited to narrowband
directional patterns: a concentric ring spreads its energy over all
orientations and is therefore outside the model, which we state and test
rather than assume.

\subsection{Task Based Assessment}

Pixel fidelity metrics are known to be imperfect proxies for diagnostic
value, and the model observer framework \cite{barrett2004} provides the
alternative. The channelised Hotelling observer with Laguerre-Gauss channels
\cite{gallas2003} is the standard instrument for signal known exactly
detection tasks and is used here with a split half estimator so that the
reported detectability is free of the optimistic bias of resubstitution. Its
inclusion changes the conclusions of the study, as Section \ref{sec:results}
reports.
